\documentclass[UTF8]{ctexart}

\usepackage{geometry}  

\usepackage{amsmath}   

\title{初中方程全解}

\author{王乐知}

\begin{document}

\maketitle

方程是数学中一个重要的数学模型,是通过构建不同量之间的等量关系来求出未知量的一种数学手段,让我们一起来学习它.

\section{方程的定义}

含有未知数的等式被叫做方程,使方程左右两边相等的值叫做方程的解(也可以被叫做方程的根),求方程的解的过程叫做解方程.

\section{解方程的依据}

解方程的依据是等式的基本性质,等式的基本性质如下:

等式的基本性质1:等式两边同时加上(或减去)一个数或式,所得结果仍是等式.

等式的基本性质2:等式两边同时乘(或除以)一个不为0的数或式,所得结果仍是等式.

注意,基本性质1没有使用前提,但是基本性质2有不为0的前提,如果不能保证不为0,则会出现增根(详见分式方程一节)。

\section{解方程的流程}

每种方程都有对应的标准形式，我们要把不规则的方程化为标准方程，再代入标准方程求解。

\section{一元一次方程}

一元一次方程是最基本的方程，其只有一个未知数，且未知数的次数是1，标准形式如下：

\[
ax+b=0
\]

解一元一次方程的流程：去分母→去括号→移项→合并同类项→系数化1

\section{二元一次方程}

\end{document}